\documentclass[10pt]{article}
\usepackage[T1]{fontenc}
\usepackage[utf8]{inputenc}
\usepackage[french]{babel}
\usepackage{amsmath,amssymb,mathrsfs,mathpazo}
\usepackage{geometry}
\geometry{a4paper,margin=2cm}
\usepackage{graphicx,array,tabularx,tabularray,colortbl,multirow,multicol,booktabs,enumitem,pifont,ulem,fancybox,diagbox,cancel}
\usepackage{pstricks,pstricks-add,pst-plot,pst-tree,pst-node,pst-text,pst-func,pst-math,pst-bspline,pst-3dplot,pst-eucl,pst-solides3d}
\let\pscalculate\undefined
\usepackage{pst-bezier}
\usepackage{tikz,pgfplots,tkz-tab}
\usetikzlibrary{arrows,arrows.meta,patterns,calc,positioning,trees,decorations.markings,shapes.geometric,shapes.misc,shapes.arrows,intersections,decorations.pathmorphing,decorations.pathreplacing,backgrounds,shadows,babel}
\pgfplotsset{compat=1.18}
\usepgfplotslibrary{fillbetween}
\usepackage[np]{numprint}
\usepackage{listings,makecell,pas-tableur,twoopt}
\let\barre\undefined
\usepackage[pstricks]{tablvar}
\definecolor{darkgreen}{RGB}{0,100,0}
\providecommand{\red}{\color{red}}
\providecommand{\blue}{\color{blue}}
\providecommand{\green}{\color{green}}
\providecommand{\black}{\color{black}}
\providecommand{\white}{\color{white}}
\providecommand{\gray}{\color{gray}}
\providecommand{\cyan}{\color{cyan}}
\providecommand{\magenta}{\color{magenta}}
\providecommand{\yellow}{\color{yellow}}
\providecommand{\violet}{\color{violet}}
\providecommand{\orange}{\color{orange}}
\providecommand{\euro}{\texteuro}
\providecommand{\EUR}{\texteuro}
\providecommand{\ieme}{\textsuperscript{e}}
\providecommand{\degres}{\ensuremath{^\circ}}
\providecommand{\degree}{\ensuremath{^\circ}}
\providecommand{\vv}[1]{\overrightarrow{#1}}
\providecommand{\Coord}[1]{#1}
\providecommand{\decosix}{}
\providecommand{\starredbullet}{\ensuremath{\bullet}}
\pagestyle{empty}
\setlength{\parindent}{0pt}
\usepackage{adjustbox}\DeclareUnicodeCharacter{2236}{:}\DeclareUnicodeCharacter{2190}{\ensuremath{\leftarrow}}
\providecommand{\R}{\mathbb{R}}
\providecommand{\N}{\mathbb{N}}
\providecommand{\Z}{\mathbb{Z}}
\providecommand{\Q}{\mathbb{Q}}
\providecommand{\D}{\mathbb{D}}
\providecommand{\C}{\mathbb{C}}
\providecommand{\e}{{\mathrm{e}}}
\providecommand{\vect}[1]{\overrightarrow{#1}}
\providecommand{\vectt}[1]{\overrightarrow{\mathrm{#1}}}
\providecommand{\vvt}[1]{\overrightarrow{\mathrm{#1}}}
\providecommand{\np}[2][]{{#2~#1}}
\providecommand{\vv}[1]{{\overrightarrow{#1}}}
\renewcommand{\d}{{\mathrm{d}}}
\providecommand{\degres}{{^\circ}}
\providecommand{\degre}{{^\circ}}
\providecommand{\blue}{\color{blue}}
\providecommand{\red}{\color{red}}
\providecommand{\enskip}{\quad}
\providecommand{\llbracket}{\lbrack\!\lbrack}
\providecommand{\rrbracket}{\rbrack\!\rbrack}
\providecommand{\euro}{{\text{€}}}
\providecommand{\no}{{\text{n°}}}
\providecommand{\interval}[2]{{[#1;#2]}}
\providecommand{\intervFO}[2]{{[#1;#2[}}
\providecommand{\dotfill}{\ldots}

\providecommand{\ds}{\displaystyle}
\providecommand{\textding}[1]{#1}
\providecommand{\Oij}{\(\mathrm{O};\vec{\imath},\vec{\jmath}\)}
\providecommand{\Oijk}{\(\mathrm{O};\vec{\imath},\vec{\jmath},\vec{k}\)}
\providecommand{\Ouv}{\(\mathrm{O};\vec{u},\vec{v}\)}
\providecommand{\intervFF}[2]{\ensuremath{[#1;#2]}}
\renewcommand{\euro}{\ifmmode\text{\texteuro}\else\texteuro\fi}
\DeclareUnicodeCharacter{00B0}{\ifmmode^{\circ}\else\textdegree\fi}
\DeclareUnicodeCharacter{20AC}{\ifmmode\text{\texteuro}\else\texteuro\fi}
% Énoncés : Éducation nationale. Transcription et corrigés : APMEP, Denis Vergès et contributeurs. Découpage et classement : CoopMaths, Sébastien Lozano, Éric Elter et contributeurs. Adaptation de la présentation : Maths Explorer.
% CC BY-SA : https://coopmaths.fr/www/about/
% https://coopmaths.fr/alea/static/bac/2024/tex/bac_2024_06_sujet2_metropole_devoile_3_cor.tex
\begin{document}
\noindent\textit{APMEP / CoopMaths. CC BY-SA. Présentation adaptée par Maths Explorer.}\par\medskip
\medskip

Soit $a$ un nombre réel strictement supérieur à 1.

On considère la suite $\left(u_n\right)$ définie par $u_0 = a$ et, pour tout entier naturel $n$ : \[u_{n+1} = u_n^2 - 2u_n + 2.\]

On admet que pour tout entier naturel $n$,\: $u_n > 1$.

L'objectif de cet exercice est d'étudier la suite $\left(u_n\right)$ pour différentes valeurs du nombre réel $a$.

\bigskip

\textbf{Partie A : étude de la suite \boldmath $\left(u_n\right)$ \unboldmath dans le cas }\boldmath$1 < a < 2$\unboldmath 

\medskip

\begin{enumerate}
\item 
	\begin{enumerate}
		\item 
D'après la définition : pour tout entier naturel $n$
		
$u_{n+1} = u_n^2 - 2u_n + 2 \iff u_{n+1}- 2 = u_n^2 - 2u_n  \iff u_{n+1}- 2 = u_n\left(u_n - 2\right)$.
		\item 
\emph{Méthode} 1 : d'après la définition : pour tout entier naturel $n$
		
$u_{n+1} = u_n^2 - 2u_n + 2 \iff u_{n+1} = \left(u_n - 1\right)^2 - 1 + 2 \iff u_{n+1} = \left(u_n - 1\right)^2  + 1 \iff  u_{n+1} - u_n = \left(u_n - 1\right)^2  + 1 - u_n \iff u_{n+1} - u_n = \left(u_n - 1\right)^2  - 1\left(u_n - 1\right) \iff  u_{n+1} - u_n = \left(u_n - 1\right)\left(u_n - 1 - 1)\iff \right] \iff u_{n+1} - u_n = \left(u_n - 1\right)\left(u_n - 2\right)$

\emph{Méthode} 2 : Soit $P_n = u_{n+1} - u_n = u_n^2 - 2u_n + 2 - u_n = u_n^2 - 3u_n + 2$ ; en posant $u_n = x$, on obtient 

$P_n = x^2 - 3x + 2$ : ce trinôme a deux racines (car $\Delta = 3^2 - 4 \times 2 = 1$) :

$x_1 = \dfrac{3 + 1}{2} = 2$ et $x_2 = \dfrac{3 - 1}{2} = 1$.

Donc $P_n = (x - 1)(x - 2) = \left(u_n - 1\right)\left(u_n - 2\right)$.
	\end{enumerate}
\item 
	\begin{enumerate}
		\item 
		\emph{Initialisation} : $u_0 = a$ et $1 < a < 2$, donc $u_0 < 2$ : l'inégalité est vraie au rang 0 ;

		\emph{Hérédité} : soit $n \in \N$ tel que $u_n < 2$.

$u_n < 2 \iff u_n - 2 < 0$ et on a admis que $u_n > 1 \iff u_n - 1 > 0$, donc d'après le \textbf{1 b.} 	$\left(u_n - 1\right)\left(u_n - 2\right) < 0$ c'est-à-dire $u_{n+1} - u_n  < 0\iff u_{n+1} < u_n < 2$, d'où par transitivité : $u_{n+1} < 2$ : l'inégalité est vraie au rang $n + 1$.

\emph{Conclusion} : l'inégalité est vraie au rang 0 et si elle est vraie au rang $n$, elle l'est aussi au rang $n + 1$, donc par le principe de récurrence :

pour tout naturel $n \in \N, \:\: 1< u_n < 2$.
		\item 
On a vu dans la question précédente que $1 < u_{n+1} < u_n < 2$ qui montre :
		
		$\bullet~$ que la suite $\left(u_n\right)$ est décroissante ;
		
		$\bullet~$ que la suite $\left(u_n\right)$ est minorée par 1.
		
		La suite  $\left(u_n\right)$ est monotone décroissante et minorée par le 1 : elle converge donc vers un réel $\ell \geqslant 1$.
		
		Par continuité de la fonction polynôme $x \longmapsto x^2 - 2x + 2$, la relation de récurrence donne à la limite en plus l'infini :
		
		$\ell = \ell^2 - 2\ell + 2 \iff \ell^2 - 3\ell + 2  = 0 \iff \left\{\begin{array}{l c l}
\ell - 1&=&0\\
&\text{ou}&\\
\ell - 2&=&0\\
\end{array}\right. \iff \left\{\begin{array}{l c l}
\ell &=&1\\
&\text{ou}&\\
\ell &=&2\\
\end{array}\right. $

$\ell = 2$ n'est pas possible puisque $\left(u_n\right)$ étant strictement décroissante : 

$\ell < u_0 = a < 2$, d'où $\ell < 2$, \:donc $\ell = 1$.
	\end{enumerate}
\end{enumerate}
	
	\bigskip
	
\textbf{Partie B : étude dans le cas particulier }\boldmath $a = 2$\unboldmath

\medskip

\begin{enumerate}
\item u(2,1) renvoie 2 et u(2,2) renvoie 2
















\item 



On peut conjecturer que la suite $\left(u_n\right)$ est constante : $u_n = 2$ quel que soit $n \in \N$
\end{enumerate}

\bigskip

\textbf{Partie C : étude dans le cas général}

\medskip

\begin{enumerate}
\item On considère la suite $\left(v_n\right)$ définie, pour tout entier naturel $n$, par $v_n = \ln \left(u_n - 1\right)$.
	\begin{enumerate}
		\item 
Quel que soit $n \in \N, \: v_{n+1} = \ln \left(u_{n+1} - 1\right) = \ln \left(u_n^2 - 2u_n + 2 - 1\right) = \ln \left[\left(u_n - 1\right)^2\right] = 2\ln \left(u_n - 1\right)$ car on sait que $u_n > 1 \iff u_n - 1 > 0$.

Finalement : $v_{n+1} = 2\ln \left(u_n - 1\right) = 2v_n$.

L'égalité $v_{n+1} = 2v_n$ vraie pour tout $n \in \N$ montre que la suite $\left(v_n\right)$ est une suite géométrique de raison 2 de premier terme $v_0 = \ln a - 1)$
		\item 
On sait que quel que soit $n \in \N, \: v_n = v_0 \times 2^n = 2^n\ln (a - 1)$.

On peut écrire  $v_n = \ln (a - 1)^{2^n} = \ln \left(u_n - 1\right)$ soit par croissance de la fonction logarithme népérien :

$(a - 1)^{2^n} = u_n - 1 \iff u_n = 1 + (a - 1)^{2^n} = 1 + \e^{2^n \ln (a - 1)}$.
	\end{enumerate}
\item ~

$\bullet~~$Si $1 < a < 2$, alors $0 < a - 1 < 1 \Rightarrow \ln (a - 1) < 0$ (par croissance de la fonction logarithme népérien).

On a donc $2^n \ln (a - 1) < 0$ et on sait que $\displaystyle\lim_{n \to + \infty} \e^{2^n \ln (a - 1)} = 0$, donc $\displaystyle\lim_{n \to + \infty}u_n = 1$.

$\bullet~~$Si $a = 2, \: \ln (a - 1) = \ln 1 = 0$, donc $2^n \ln (a - 1) = 0$ et $u_n = 1 + 1 = 2$. $u$ est constante et on peut écrire $\displaystyle\lim_{n \to + \infty}u_n = 2$.

$\bullet~~$ Si $a > 2, \: a - 1 > 1$, donc $\ln (a - 1) > 0$ et $\displaystyle\lim_{n \to + \infty} \e^{2^n \ln (a - 1)} = + \infty$, donc $\displaystyle\lim_{n \to + \infty}u_n = + \infty$. La suite est divergente.
\end{enumerate}

\bigskip

\end{document}
